\scanpage{52}

% B1 | text | opening lines
Á l'infini, la situation n'est pas très belle encore. Il
vaut mieux dire :

% B2 | text | first full paragraph
Pour tt $\xi_p\in\mathcal T_p$, \deleted{la valeur en $\infty$ de $\mathscr G(\xi)$ est} \inserted{considérer la gerbe $\unsure{h|\mathfrak h}(\xi_p)$} de
groupe $\mu_2$ sur $\unsure{\mathbb Z_{1/p}\mathbb Z|\mathbb Z[1/p]}$ définie par $(C(\xi_p),T(\xi_p))$, \deleted{En fait} \inserted{Considérons}
$\deleted{\unsure{\mathscr G}} h(\xi_p)_{\mathbb R}$, qui est une gerbe \uline{canonique} sur $\mathbb R$ de groupe
$\mu_2$. \deleted{J'ignore si} \inserted{pour $\xi_p=$ élément de Deuring,} elle a une description indépendante
de la théorie des catégories de Deuring (et de celle
des courbes elliptiques). \unsure{Ce qu'on|Ce que} \unsure{sait|sont} en tous cas
qu'elle est non triviale ; elle est, pour $\xi_p=$ élément de Deuring,
munie d'une polarisation $\xi_\infty$ (ce qui revient à la
donnée d'une classe d'isomorphie \inserted{d'équivalences} avec la gerbe
type $g_\infty$, \uline{mais il reste} \uline{l'ambiguïté du signe}.

% B3 | text | middle, category over the reals
Catégorie $C(\xi_p)_\infty$ \deleted{\unsure{est}} $\simeq$ catégorie des vectoriels \inserted{à droite $V$ de dim $1$} \deleted{\unsure{qui}} sur
$\mathbb H$ ($=$ quaternions ordinaires), avec
\deleted{\unsure{leur} $V$} foncteur norme
réduite vers cat. de Picard de $\mathbb R$.

% B4 | text | lower middle, functor
Ainsi il y a un foncteur ``canonique au signe près''

% B5 | diagram | lower middle, arrow and lower additions
\diagram{\begin{tikzcd}[ampersand replacement=\&,column sep=large]
C(\xi_p)\arrow[r,"{\deleted{\unsure{\xi}}c_\infty}"]\&\deleted{\unsure{C(\mathbb R)}\ \text{catégorie}}\,C(\mathbb R)\\
(T(\xi_p)\arrow[r,draw=none]\&\deleted{\unsure{\xi}}\;\mathbb R)
\end{tikzcd}}

% B6 | text | left-hand leader and right-hand gloss
\margin{$\longleftarrow$ n'est pas spécial à $C(\xi_p)$, vaut \deleted{\unsure{pour}} pour tt $C$}
\margin{la \unsure{factorisation|fonction} \unsure{Arr.} par les vectls quaternioniens de dim $1$ \uuline{unitaires}}

% B7 | text | struck lines below diagram, with continuation
\deleted{et qui associe à chaque élément de
\unsure{la} \unsure{leur} \unsure{définit} un autre}
tout autre choix étant isom
à ça, l'isom \uline{défini au signe près} (\unsure{chic|chicos} !).

% B8 | text | foot, last paragraph
Prenant $\xi_p=$ élément de \deleted{\unsure{Eisenstein}}\inserted{Deuring}, on obtient un
foncteur \inserted{ess. canonique} des courbes elliptiques de Hasse nul
vers les vectoriels \inserted{\unsure{quatrs} unitaires} de dim $1$. La \deleted{\unsure{transp}} Norme
correspond à la polarisation, la \unsure{transposition|transpiration}, la transposition
% archive sheet number: 51
